\clearpage\thispagestyle{empty}
\pdfbookmark[0]{Quelle zum Brief an von Lindenau vom 6. August 1808}{lindenau1808aug06-sources}
\begin{center}\textbf{Quelle zum Brief an von Lindenau vom 6. August 1808}\end{center}
\noindent Diese KI-Übertragung ins heutige Deutsch folgt dem Brief von Carl Friedrich Gauß an Bernhard August von Lindenau vom 6. August 1808 in Gauß’ \href{https://gdz.sub.uni-goettingen.de/id/PPN236007467}{\textit{Werke}, Bd.~VI} (Göttingen, 1874), S.~305--306. Der Erstdruck erschien als Nr.~XXI, \textit{Auszug aus einem Schreiben des Hrn. Dr. Gauß}, im Augustheft der \textit{Monatlichen Correspondenz}, Bd.~XVIII, S.~\href{https://www.digitale-sammlungen.de/en/view/bsb10538609?page=193}{173}, \href{https://www.digitale-sammlungen.de/en/view/bsb10538609?page=194}{174} und \href{https://www.digitale-sammlungen.de/en/view/bsb10538609?page=195}{175}.
\par\medskip
Die Werkausgabe enthält den handschriftlichen Eintrag zu Ceres und verweist für die Vestabeobachtungen auf die Mitteilung vom November. Für den Satz über die noch in Paris befindlichen Instrumente verweist Ernst Schering auf die Göttingischen gelehrten Anzeigen, Stück 127, S.~\href{https://manifests.sub.uni-goettingen.de/iiif/image/gdz:PPN319721507_1808:00001271/full/1600,/0/default.jpg}{1265}. Mit \textit{levi calamo} sagt Gauß, dass er die Rechnung rasch und ohne Umstände angestellt hat. Zu den Instrumentenbezeichnungen siehe den \href{https://kalliope.staatsbibliothek-berlin.de/ead?ead.id=DE-611-HS-3300344}{Katalogeintrag der SUB Göttingen zu Cod. Ms. Gauß Astron.~7}. KI-Übertragung und KI-Satz.
\clearpage
